\documentclass[10pt,letterpaper]{article}
\usepackage{cornell-notes}

\title{Clause 27.08.07.04: Set Intersection}
\author{}
\date{}
\setDocTitle{Clause 27.08.07.04: Set Intersection}
\setDocSubtitle{C++ 2024}
\setDocOwner{C++ 2024 Cornell Notes}
\setDocDate{\today}
\setCornellCollection{C++ 2024 Cornell Notes}
\setCornellUnitType{Clause}
\setCornellUnitNumber{27.08.07.04}
\setCornellUnitTitle{Set Intersection}
\hypersetup{pdftitle={Clause 27.08.07.04: Set Intersection},pdfsubject={Cornell notes},pdfauthor={}}

\begin{document}
\maketitle
\begin{CornellOverviewBox}{Overview}
\textbf{Classification:} Standard library - Normative\\[2pt]
\textbf{Learning objective:} Understand the rules, terminology, and semantic consequences associated with set\_intersection.
\end{CornellOverviewBox}\begin{CornellSummaryBox}{Summary}
{\sffamily\bfseries\color{Accent} Summary}\par\vspace{3pt}Section 27.8.7.4 focuses on set\_intersection. Apply the specified interface together with its mandates, effects, result conditions, exception behavior, and complexity constraints. When applying it, preserve the distinction between mandatory requirements, recommendations, permissions, and behavior deliberately left undefined, unspecified, or implementation-defined.
\end{CornellSummaryBox}

\end{document}
